\documentclass[../principal.tex]{subfiles}
\graphicspath{{\subfix{../imagenes/}}}

\begin{document}

\chapter{Color}

En este capítulo veremos definiciones de color, la percepción humana de color, el modelamiento de color usando computadores, y su aplicación en p5.js

\section{Definición}

Definiremos color como la percepción de luz a través de nuestros ojos.

La velocidad de la luz es de trescientos mil kilómetros por segundo, en números  $300.000 \frac{km}{s}$.

\section{Percepción humana de color}

Se estima que una persona es capaz de detectar alrededor de 10 millones de colores distintos.

La percepción ocurre entre tonalidades rojas, aproximadamente con una frecuencia de $400 THz$, y tonalidades violeta, con una frecuencia aproximada de $800 THz$.

Las ondas de luz que tienen una frecuencia menor a rojo no son visibles y son denominadas infrarrojas, como las que se usan para la emisión de señales de controles remotos domésticos.

Las ondas de luz que tienen una frecuencia mayor a violeta no son visibles y son denominadas ultravioleta, como las que producen quemaduras en la piel.

La percepción en nuestros ojos ocurre con receptores de 3 rangos de frecuencia distintos, en torno a los colores rojo, verde y azul.

\section{Color y computación}

Para describir un color necesitamos un sistema que sea capaz de tener al menos 10 millones de combinaciones, que es la cantidad de colores que una persona puede percibir.

1 bit puede describir 2 valores posibles: 0 y 1. Si tenemos 2 bits, tenemos 4 valores posibles: 00, 01, 10 y 11. Con 3 bits, tenemos 8 valores posibles: 000, 001, 010, 011, 100, 101, 110, 111. La ecuación que gobierna cuántas combinaciones distintastenemos con N bits es:

$${combinaciones} = {2}^{N}$$

Entonces lo que estamos buscando, es cuántos bits necesitamos para tener al menos diez millones de valores posibles.

Si aumentamos N, podemos calcular que con N=23 y N=24:

$${2}^{23} = 8388608$$
$${2}^{24} = 16777216$$

Entonces con N = 23 no nos alcanza, recién con N = 24 superamos los 10 millones, y ese es el valor elegido por la industria computacional para representar colores.

Debido a la percepción humana de color, estos 24 bits se usan para describir cuán encendida están 3 ampolletas o canales de colores: una dedicada a rojo, una a verde, y una azul. Este sistema se denomina RGB por sus iniciales en inglés.

Cada canal de color entonces recibe 8 bits de resolución, y con eso cada canal tiene $2^8 = 256$ valores posibles. Y como en computación contamos desde 0, los valores de cada canal de color están entre el rango 0 y 255, para representar todo el espectro entre apagado y prendido.

\section{Representación hexadecimal}

Como cada canal tiene 8 bits de resolución, para describir un color necesitamos los 3 canales y con esto, el color rojo por ejemplo recibe la siguiente notación:

$1111 1111 0000 0000 0000 0000$

Si tomamos los 8 bits de cada canal, y los dividimos entre los primeros 4 y los siguientes 4, entonces tenemos que los 24 bits son 6 grupos de 4 bits.

Los 4 bits permiten representar los siguientes números en binario:

\begin{table}
\caption{Representaciones binarias, decimales y hexadecimales de valores entre 0 y 15}
\label{tab:starch}
\tbfont
\begin{tabular}{@{}lll}
\toprule
\thfont Binario & \thfont Decimal & \thfont Hexadecimal \% \\
\midrule
0000 & 0  & 0 \\
0001 & 1  & 1 \\
0010 & 2  & 2 \\
0011 & 3  & 3 \\
0100 & 4  & 4 \\
0101 & 5  & 5 \\
0110 & 6  & 6 \\
0111 & 7  & 7 \\
1000 & 8  & 8 \\
1001 & 9  & 9 \\
1010 & 10 & A \\
1011 & 11 & B \\
1100 & 12 & C \\
1101 & 13 & D \\
1110 & 14 & E \\
1111 & 15 & F \\
\bottomrule
\end{tabular}
\end{table}

\section{Referencias}

\begin{enumerate}
    \item \url{https://en.wikipedia.org/wiki/Color}
\end{enumerate}


\end{document}
